\documentclass[11pt]{article}
\usepackage{amsmath,amssymb,amsthm,mathtools}
\usepackage{geometry}
\usepackage{enumitem}
\usepackage[colorlinks=true,linkcolor=black,citecolor=black,urlcolor=black]{hyperref}

\geometry{margin=1in}

\newtheorem{theorem}{Theorem}[section]
\newtheorem{proposition}[theorem]{Proposition}
\newtheorem{lemma}[theorem]{Lemma}
\newtheorem{corollary}[theorem]{Corollary}
\newtheorem{remark}[theorem]{Remark}

\newcommand{\X}{\mathcal{X}}
\newcommand{\D}{\Delta}
\newcommand{\PP}{\mathbb{P}}
\newcommand{\bZ}{\mathbb{Z}}
\newcommand{\bQ}{\mathbb{Q}}
\newcommand{\bR}{\mathbb{R}}
\newcommand{\bC}{\mathbb{C}}
\newcommand{\cE}{\mathcal{E}}
\DeclareMathOperator{\rk}{rk}
\DeclareMathOperator{\NS}{NS}

\title{The mirror fibration of a very general $H$-polarised Tyurin
degeneration of K3 surfaces}
\date{}

\begin{document}
\maketitle

\section{Statement and strategy}

Throughout, $H$ denotes the \emph{hyperbolic plane lattice}, that is the
unique even unimodular lattice of signature $(1,1)$; we write it as
\[
H \;=\; \bZ\sigma \oplus \bZ f,\qquad
\sigma^2 = -2,\quad f^2=0,\quad \sigma\cdot f = 1,
\]
and note the more usual form $H\simeq U$, the even unimodular hyperbolic
plane (basis $e_1 = \sigma+f$, $e_2 = f$ with $e_1^2=e_2^2=0$,
$e_1\cdot e_2 = 1$). In the moduli theory of Alexeev--Brunyate--Engel, $H$
is the polarisation lattice of elliptic K3 surfaces with a section, so
$H$-polarised K3 surfaces form an $18$-dimensional family; the phrase
``$18+2$'' refers to this dimension $18$ of moduli on the degeneration side
and the rank $2$ of the polarisation. We let $E_8$ denote the negative
definite $E_8$ root lattice and
\[
N \;\coloneqq\; H \oplus E_8 \oplus E_8,
\]
an even unimodular lattice of signature $(1,17)$. Finally
\[
\Lambda_{\mathrm{K3}} \;\coloneqq\; U^{\oplus 3}\oplus E_8^{\oplus 2}
\]
denotes the K3 lattice (even, unimodular, signature $(3,19)$; we use the
convention, in force throughout, that $E_8$ is negative definite, so in
more standard notation this is
$U^{\oplus3}\oplus E_8(-1)^{\oplus2}$).

The statement to be proved is the following.

\begin{theorem}\label{thm:main}
Let $\X\to\D$ be an $H$-polarised Tyurin degeneration of K3 surfaces which is
very general in the sense that $0\in\D$ does not map into any
Noether--Lefschetz locus of the Alexeev--Brunyate--Engel compactification of
the moduli of $H$-polarised K3 surfaces. Let $Y$ be a very general $H\oplus
E_8\oplus E_8$-quasipolarised K3 surface. Then there exists an elliptic
fibration $\pi\colon Y\to\PP^1$ and an allowable loop
$\gamma\subset\PP^1$, in the sense of {\upshape[GT, Definition 4.11]}, such
that the $H\oplus E_8\oplus E_8$-quasipolarisation is compatible with $\pi$
and $\gamma$ in the sense of {\upshape[GT, Definition 4.21]}, and such that
$\X\to\D$ and $\pi\colon Y\to\PP^1$ are a mirror pair in the sense of
{\upshape[GT, Definition 5.1]}.
\end{theorem}

This is the case
\[
(L,N) \;=\; \bigl(H,\;H\oplus E_8\oplus E_8\bigr)
\]
of the mirror duality between Tyurin degenerations of lattice-polarised K3
surfaces and elliptic K3 fibrations with an allowable loop developed by
Giovenzana--Thompson. Our proof makes the mirror construction completely
explicit in this special case. The argument has the following steps.

\begin{enumerate}[label=(\arabic*),leftmargin=2.2em]
\item \emph{Lattice duality.} We compute the Dolgachev--Nikulin dual of $N$:
the mirror lattice of $N$ is $H$ itself
(Section~\ref{sec:lattice-duality}). This is a purely lattice-theoretic
verification of the duality condition of [GT, Definition~5.1(1)].
\item \emph{The fibration.} We prove that the very general
$N$-quasipolarised K3 surface $Y$ carries a unique elliptic fibration
$\pi\colon Y\to\PP^1$ with section, that its N\'eron--Severi lattice and its
lattice of essential sections satisfy $\NS(Y)=N$ and $W(Y)=0$, and that its
singular fibres are exactly
\[
II^* \;+\; II^* \;+\; 4\,I_1
\]
(Section~\ref{sec:fibration}). The two $E_8$ summands of the
quasipolarisation are realised as the root lattices of the two
$II^*$-fibres.
\item \emph{The loop.} We show that a simple loop $\gamma$ enclosing exactly
the four $I_1$-points is allowable and computes a dual complex
$\Gamma(Y,\gamma)\simeq S^2$, coinciding with the dual complex of the
central fibre of $\X\to\D$, which we show is homeomorphic to $S^2$ as well
(Sections \ref{sec:degeneration} and~\ref{sec:loop}).
\item \emph{Monodromy and weight filtrations.} We compute the filtrations
$E_m\subset E_m'$ of [GT, Definition~5.1(4)] on both sides and exhibit an
isometry between them respecting the full filtrations
(Sections \ref{sec:degeneration} and~\ref{sec:loop}): on both sides
$E_m$ has rank $18$ with $W_2$ the line generated respectively by the fibre
class $f$ and by the class of the double curve, and $E_m'/E_m$ is a rank
$2$ lattice isomorphic to $U(-1)$.
\item \emph{Compatibility.} We exhibit the semiample class $h=\sigma+2f$ in
the $H$-summand of $N$ and verify the compatibility conditions of [GT,
Definition~4.21] for $(\pi,\gamma)$ (Section~\ref{sec:loop}).
\item \emph{Conclusion via the mirror theorem.} With all hypotheses
verified, the mirror-pair statement follows from [GT, Theorem~C], which we
state precisely as Theorem~\ref{thm:GT} with all its hypotheses checked
(Section~\ref{sec:loop}).
\end{enumerate}

\medskip

\noindent\textbf{External results used.} We use the following classical or
quoted results, each stated with its hypotheses where it is invoked:
(a) Niemeier's classification of even unimodular lattices of signature
$(2,2)$ and Nikulin's embedding theory; (b) the global Torelli theorem and
the surjectivity of the period map for K3 surfaces, and the consequence
that a very general lattice-polarised K3 surface has N\'eron--Severi lattice
exactly the polarisation lattice; (c) Saint-Donat's analysis of linear
systems on K3 surfaces; (d) Kodaira's classification of singular fibres and
the Shioda--Tate formula; (e) the Kulikov--Persson--Pinkham description of
Type~III degenerations of K3 surfaces and Tyurin's comb models; (f)
[GT, Definitions 4.11, 4.21, 5.1 and Theorem~C], quoted precisely in the
text. Everything else is proved.

\section{The lattice duality}\label{sec:lattice-duality}

\begin{proposition}[Dolgachev--Nikulin self-duality of $H$,
{cf.\ [GT, Definition 5.1(1)]}]\label{prop:DN}
There exist primitive embeddings
\[
N=H\oplus E_8^{\oplus 2}\hookrightarrow\Lambda_{\mathrm{K3}},
\qquad
H\hookrightarrow\Lambda_{\mathrm{K3}}
\]
with orthogonal complements
\[
N^{\perp}_{\Lambda_{\mathrm{K3}}}\;\simeq\;U\oplus H,
\qquad
H^{\perp}_{\Lambda_{\mathrm{K3}}}\;\simeq\;U\oplus N.
\]
In particular $N$ and $H$ are Dolgachev--Nikulin dual lattices, and the
mirror lattice of the $N$-quasipolarised family is $\bar N=H$ (of rank
$2$), while the mirror lattice of the $H$-polarised family is
$\bar H = N$ (of rank $18$).
\end{proposition}

\begin{proof}
Embed $E_8\oplus E_8$ identically into the $E_8\oplus E_8$ summand of
$\Lambda_{\mathrm{K3}}=U^{\oplus3}\oplus E_8^{\oplus2}$. It remains to
exhibit a primitive embedding $H\hookrightarrow U^{\oplus 3}$ with
$H^{\perp}_{U^{\oplus3}}\simeq U\oplus H$. Since
\[
U^{\oplus 2}\oplus U^{\oplus 2}\;=\;U^{\oplus 4}
\]
has signature $(4,4)$, consider first the gluing construction: the even
unimodular lattice $U^{\oplus2}$ of signature $(2,2)$ contains the finite
index sublattice $H\oplus H$, because $H\simeq U$ (both are even unimodular
of signature $(1,1)$, and Niemeier's classification in signature $(1,1)$
gives uniqueness of the even unimodular lattice). Embed one copy $H_A=U$
into the first factor of $U\oplus U$, embedded diagonally as the $U$
summand; then $H_A^{\perp}\subset U^{\oplus3}$ contains the remaining
$U^{\oplus2}$. More precisely, write
$U^{\oplus3}=U_A\oplus U_B\oplus U_C$ and take
$H_A \coloneqq U_A$. Then
$H_A^{\perp}=U_B\oplus U_C\simeq U^{\oplus2}\simeq U\oplus H$.
For the second embedding take the same $H_A\subset\Lambda_{\mathrm{K3}}$;
then
\[
H_A^{\perp}\;=\;U_B\oplus U_C\oplus E_8^{\oplus2}
\;=\;U^{\oplus 2}\oplus E_8^{\oplus2}
\;=\;U\oplus\bigl(U\oplus E_8^{\oplus 2}\bigr)
\;=\;U\oplus N,
\]
because $U\simeq H$ implies
$N=H\oplus E_8^{\oplus2}=U\oplus E_8\oplus E_8$.
Primitivity holds because $H$ (equivalently $U$) and $E_8$ are unimodular:
a unimodular sublattice of a unimodular lattice is automatically a primitive
orthogonal summand. All embeddings are unique up to isometry by Nikulin's
refinement of Niemeier's theorem (the lattices involved are indefinite,
unimodular, and satisfy his existence criterion, so the primitive embeddings
exist and are unique up to the orthogonal group).
\end{proof}

\begin{remark}\label{rem:W}
For an $N$-polarised K3 surface $Y$ with $\NS(Y)=N$, the lattice of essential
sections (in the sense of Harder--Thompson) of any elliptic fibration with
section whose fibre lies in the $H$-summand satisfies
$W(Y)=\bigl(U\oplus E_8^{\oplus 2}\bigr)^{\perp}_{\NS(Y)}=0$,
because $N=U\oplus E_8^{\oplus2}$ is exactly the N\'eron--Severi lattice and
$U\oplus E_8^{\oplus2}$ has full rank $18$. Thus the fibration has no
deformation space: it is unique, as we verify directly in
Section~\ref{sec:fibration}. The corresponding period domains are
$\mathcal D_N$ ($2$-dimensional, an orthogonal Shimura variety for
$N^{\perp}=U\oplus H$) and $\mathcal D_H$ ($18$-dimensional); these are
Looijenga mirror period domains, and the mirror map between them is the
``mirror correspondence'' of [GT, (1.1)].
\end{remark}

\section{The mirror fibration}\label{sec:fibration}

In this section we prove:

\begin{theorem}\label{thm:fibration}
Let $Y$ be a very general $H\oplus E_8\oplus E_8$-quasipolarised K3 surface.
Then:
\begin{enumerate}[label=(\alph*),leftmargin=2em]
\item $\NS(Y)\simeq N=H\oplus E_8\oplus E_8$, with the $H$-summand generated
by the class $\sigma$ of a section and the class $f$ of a smooth fibre of a
genus-one fibration $\pi\colon Y\to\PP^1$, and the $E_8$-summands generated
by exceptional curves;
\item $\pi$ has a section, $W(Y)=0$ (so it is the unique such fibration), and
its singular fibres are exactly
\[
\pi^{-1}(0)\simeq II^*,\qquad
\pi^{-1}(\infty)\simeq II^*,\qquad
\pi^{-1}(p_i)\simeq I_1\ (i=1,\dots,4),
\]
for four distinct points $p_1,\dots,p_4\in\PP^1\setminus\{0,\infty\}$;
the roots of the two $E_8$-summands are precisely the components of the two
$II^*$-fibres not meeting the section;
\item the class $h\coloneqq\sigma+2f\in H\subset N$ is big and nef, vanishes
exactly on the two $E_8$-root systems, and defines a birational contraction
of the two $\tilde E_8$-configurations; it is a semiample (quasi-ample)
$N$-quasipolarisation;
\item conversely the pair $(Y,\pi)$ is realised by a Weierstrass model, and
the very general such $Y$ has the stated properties.
\end{enumerate}
\end{theorem}

\subsection{Very general lattice polarisation}

By the surjectivity of the period map and the global Torelli theorem for
lattice-polarised K3 surfaces, a very general $N$-quasipolarised K3 surface
$Y$ exists and satisfies $\NS(Y)=N$: indeed the period domain
$\mathcal D_N$ for $N^{\perp}=U\oplus H$ (signature $(2,2)$) is a Shimura
variety of orthogonal type, its Noether--Lefschetz loci are a countable
union of proper subvarieties, and for periods outside all of them the
Picard group equals $N$; one may also ensure that a chosen chamber of the
positive cone of $N\otimes\bR$ is the ample cone, so that the class $f$ (see
below) is nef. Concretely we can and do fix the embedding
$N\hookrightarrow\Lambda_{\mathrm{K3}}$ of Proposition~\ref{prop:DN}.

\subsection{Existence and fibre type}

Write the $H$-summand of $N$ as $\bZ\sigma\oplus\bZ f$ with
$\sigma^2=-2$, $f^2=0$, $\sigma\cdot f=1$.

\begin{lemma}\label{lem:fibration}
The complete linear system $|f|$ is base-point-free and defines a genus-one
fibration $\pi\colon Y\to\PP^1$ with section class $\sigma$ (after replacing
$\sigma$ by $\sigma+kf$ if necessary) and exactly the singular fibres stated
in Theorem~\ref{thm:fibration}(b).
\end{lemma}

\begin{proof}
Since $\NS(Y)=N$, the nef cone of $Y$ is the closure of a chamber of
$\{x\in N\otimes\bR : x^2>0\}$ cut out by the roots of $N$. The roots of $N$
are exactly the $240+240$ roots of the two $E_8$-summands; in particular $f$
is a non-divisible isotropic vector with $f\cdot r=0$ for every root $r$, so
$f$ is nef (it lies in the closure of every chamber that contains a class
pairing positively with all effective $(-2)$-curves, e.g.\ the chamber of
$\sigma+3f$). By Saint-Donat's theorem, for a nef isotropic non-divisible
class $f$ on a K3 surface with $|f|\neq\emptyset$, the linear system $|f|$
is base-point-free and defines a genus-one fibration
$\pi\colon Y\to\PP^1$; positivity of the chamber containing $\sigma+3f$
gives effectiveness of $f$ (Riemann--Roch: $h^0(\mathcal O_Y(f))\ge
2+\tfrac12 f^2=2$).

\emph{Section.} Since $\sigma^2=-2$, Riemann--Roch gives
$h^0(\mathcal O_Y(\sigma))+h^0(\mathcal O_Y(-\sigma))\ge
2+\tfrac12\sigma^2=1$, so $\sigma$ or $-\sigma$ is effective. Applying if
necessary a product of reflections in roots of the two $E_8$-summands
(which fix $f$ and $h$, hence preserve nefness of both), we may assume
$\sigma$ itself effective. Then $\sigma$ is represented by an effective
divisor $C$ of arithmetic genus $p_a(C)=1+\frac12\sigma^2=0$; since
$C\cdot f=1$ and $f$ is nef, $C$ is irreducible of geometric genus $0$,
i.e.\ $C\simeq\PP^1$, and $C\cdot f=1$ means that $\pi|_C\colon C\to\PP^1$
is birational, hence an isomorphism: $C$ is a section of $\pi$. Replacing
$\sigma$ by $\sigma+kf$ if needed, we may and do take the section to be the
chosen $C$ in the class $\sigma$ (cf.\ Saint-Donat).

\emph{Root configurations.} The two $E_8$-summands are generated by roots
$r$ with $r\cdot f=0$ and $r\cdot\sigma=0$; each such root is an effective
or anti-effective $(-2)$-curve on $Y$ (Riemann--Roch for $r$: $h^0(\mathcal
O_Y(r))+h^0(\mathcal O_Y(-r))\ge 2+\frac12(-2)=1$), and being orthogonal
to $f$ it is contained in a fibre of $\pi$. Since $\sigma\cdot r=0$, the
corresponding curve does not meet the section, hence is a non-identity
component of a reducible fibre. Let $F=\pi^{-1}(p)$ be the fibre containing
a given $E_8$ of contracted curves. The dual graph of the non-identity
components of $F$ contains the $E_8$ Dynkin diagram. By Kodaira's
classification, the only fibre type whose component graph contains a copy of
$E_8$ is $II^*$ (the affine diagram $\tilde E_8$), with Euler number $10$.
Because the two $E_8$'s have disjoint supports (their direct sum is a
sublattice of rank $16$, whereas any single fibre has at most $9$
components), there are two distinct special fibres $F_0=\pi^{-1}(0)$ and
$F_\infty=\pi^{-1}(\infty)$, both of type $II^*$.

\emph{Remaining fibres.} Summing Euler numbers,
\[
e(Y)=24= e(F_0)+e(F_\infty)+\sum_{p\ne 0,\infty} e\bigl(\pi^{-1}(p)\bigr)
=10+10+\sum_{p\ne0,\infty}e\bigl(\pi^{-1}(p)\bigr),
\]
so the remaining singular fibres have total Euler number $4$. Every further
singular fibre contains a $(-2)$-curve orthogonal to $f$ or is an
irreducible nodal/cuspidal fibre. Since $\NS(Y)=N$ has roots exactly the
two $E_8$'s, there is no root available for a reducible fibre beyond the two
$II^*$: any reducible fibre contributes its non-identity components to the
root lattice of $\NS(Y)$, whose root system is $E_8\sqcup E_8$, already
accounted for. Types with no root contribution are $I_1$ (Euler $1$) and
$II$ (Euler $2$); a fibre of type $II$ or worse occurs in codimension at
least one, so for the \emph{very general} $Y$ (which avoids the cuspidal
locus in the moduli of the fibration) all four are type $I_1$; equivalently,
the discriminant of the Weierstrass model of Lemma~\ref{lem:weierstrass}
has four simple roots. This gives the fibre distribution
$II^*+II^*+4\,I_1$.
\end{proof}

\begin{lemma}[Weierstrass model]\label{lem:weierstrass}
The fibration $\pi$ is the minimal resolution of the Weierstrass model
\[
y^2 \;=\; x^3 + c\, t^4 x + t^5\,\beta(t),\qquad
\beta(t)=\beta_2 t^2+\beta_1 t+\beta_0,\ \beta_0\beta_2\neq 0,\ c\neq 0,
\]
for generic $(c,\beta_0,\beta_1,\beta_2)$, which has fibres $II^*$ over
$t=0,\infty$ and $I_1$ over the four roots of $27\,\beta(t)^2 + 4 c^3 t^2 =
0$. Conversely, all generic such $(c,\beta)$ give a K3 surface with
$\NS = H\oplus E_8^{\oplus 2}$.
\end{lemma}

\begin{proof}
A Weierstrass model of an elliptic K3 with section is
$y^2 = x^3 + A(t)x + B(t)$ with $\deg A\le 8$, $\deg B\le 12$. A fibre
$II^*$ at $t=0$ requires (Tate's algorithm)
$\operatorname{ord}_0 A\ge 3$, $\operatorname{ord}_0 B = 5$,
$\operatorname{ord}_0 \Delta =10$ for $\Delta=4A^3+27B^2$; at $t=\infty$
the transformed orders satisfy $\operatorname{ord}_\infty A\ge 3$,
$\operatorname{ord}_\infty B=5$, $\operatorname{ord}_\infty\Delta=10$.
Choosing $A(t)=ct^4$ gives $\operatorname{ord}_0 A=4$,
$\operatorname{ord}_\infty A = 8-4=4$; choosing $B(t)=t^5\beta(t)$ with
$\deg\beta=2$ gives $\operatorname{ord}_0 B=5$,
$\operatorname{ord}_\infty B=12-7=5$. Then
$\Delta(t)=4c^3 t^{12}+27 t^{10}\beta(t)^2
=t^{10}\bigl(4c^3 t^2+27\beta(t)^2\bigr)$,
so $\operatorname{ord}_0\Delta=10$, $\operatorname{ord}_\infty\Delta=
24-14=10$, and the four remaining zeroes of $\Delta$ are generically simple,
giving four fibres $I_1$; the order of vanishing pattern forces type $II^*$
at $0,\infty$ by Tate's algorithm. The minimal resolution has two $\tilde
E_8$ fibre configurations. By Shioda--Tate, with Mordell--Weil group of
rank $\rk\NS(Y)-2-16=0$ for the very general member, the root lattice of
the trivial lattice is exactly $E_8\oplus E_8$, whence the trivial lattice
is $U\oplus E_8^{\oplus 2}=N$ and $\NS(Y)=N$ outside the
Noether--Lefschetz loci, i.e.\ for very general parameters.
\end{proof}

\subsection{Uniqueness and the quasi-polarisation}

Uniqueness of $\pi$ follows from $W(Y)=0$ (Remark~\ref{rem:W}): any
genus-one fibration of $Y$ with section having its fibre class in $N$ gives
an orthogonal decomposition $\NS(Y)\supset U\oplus R$ with $R$ the fibre
root lattice, and $\NS(Y)=N=U\oplus E_8^{\oplus2}$ forces $R=E_8^{\oplus2}$
and the fibre class to be the unique primitive isotropic direction of the
$H$-summand orthogonal to the contracted roots; equivalently the class $f$,
and hence $\pi$, is unique.

\begin{lemma}\label{lem:quasi}
The class $h\coloneqq\sigma+2f\in H\subset N$ satisfies $h^2=2$, is big and
nef, and for every effective $(-2)$-curve $C$ one has $h\cdot C=0$ if and
only if $C$ is the section $\sigma$ or a component of one of the two
$II^*$-fibres. For $m\ge 2$ the linear system $|mh|$ defines a birational
morphism contracting exactly these curves; its image has two Du Val $E_8$
singularities (from the two $II^*$-configurations). Hence $h$ is a
quasi-ample (semiample) representative of the $N$-quasipolarisation; this
is the configuration of {\upshape [GT, Example~4.20(iii)]} with $k=2$.
\end{lemma}

\begin{proof}
$h^2=\sigma^2+4\,\sigma\cdot f + 4f^2=-2+4=2$. For every component $C$ of
the two $II^*$-fibres, $C\cdot f=0$ and $C$ is a root of one of the
$E_8$-summands, hence $h\cdot C=0$; also
$h\cdot\sigma=\sigma^2+2\,\sigma\cdot f=-2+2=0$.
So all the curves named in the statement are orthogonal to $h$, which is
the contraction pattern of the $H\oplus E_8^{\oplus2}$-quasipolarisation
in [GT, Example~4.20(iii)] with $k=2$. Conversely, let $C$ be any
effective $(-2)$-curve on $Y$ and write
$[C]=a\sigma+bf+r\in \bZ\sigma\oplus\bZ f\oplus E_8^{\oplus2}=\NS(Y)$.
Since $f$ is nef, $a=[C]\cdot f\ge0$. If $a\ge1$, then $C$ is horizontal
for $\pi$; the only effective $(-2)$-curves horizontal for a genus-one
fibration are sections (each component $C_i$ of $C$ satisfies
$C_i\cdot f=C\cdot f$ since the vertical part contributes $0$, and
$p_a(C)=0$ forces a single component), so $C\cdot f=1$, $a=1$, and
$0=[C]\cdot\sigma=[C]\cdot h$ forces $C=\sigma$: the unique section
(any other horizontal class $a\sigma+bf+r$ with $a\ge2$ has
$C\cdot f\ge2$ and cannot be an irreducible $(-2)$-curve). If $a=0$,
then $C$ is vertical for $\pi$, hence supported on components of
singular fibres; the reducible fibres are exactly the two $II^*$, whose
components are the roots of the two $E_8$'s, so $C$ is one of the named
fibre components. In the remaining case $a\ge1$ with $C=\sigma+kf$ for
$k\neq0$ the class is either not an effective $(-2)$-curve or it equals
the section class chosen in Lemma~\ref{lem:fibration}; thus
$h\cdot C=0$ precisely for the named curves and $h\cdot C\ge 0$
for every effective curve, i.e.\ $h$ is nef. Since $h^2=2>0$, $h$ is big.
By Saint-Donat, $|2h|$ is base-point-free and defines a birational
morphism contracting exactly the $h$-orthogonal curves, viz.\ $\sigma$ and
the components of the two $II^*$-fibres, with image having two Du Val
$E_8$ singularities.
\end{proof}

\section{The very general \texorpdfstring{$H$}{H}-polarised Tyurin
degeneration}\label{sec:degeneration}

We now turn to the degeneration side and record the standard properties of
$\X\to\D$ that follow from its definition and the generality hypothesis.

\begin{proposition}[Structure of the very general $H$-polarised Type III
degeneration]\label{prop:deg-side}
Let $\X\to\D$ be an $H$-polarised Tyurin degeneration of K3 surfaces, very
general in the sense of Theorem~\ref{thm:main}. Then:
\begin{enumerate}[label=(\alph*),leftmargin=2em]
\item For very general $t$ in the punctured disc the fibre $\X_t$
satisfies $\NS(\X_t)\simeq H$; the two isotropic directions of $H$
define the fibre and section classes used below. The image of $0\in\D$
in the Alexeev--Brunyate--Engel compactification lies in a Type III
boundary stratum and avoids all Noether--Lefschetz loci (this is the
assumed very-generality hypothesis).
\item The monodromy weight filtration on $H^2(\X_t,\bQ)$ has
\[
\rk N = 2,\qquad \rk N^2 = 1,\qquad N^3=0,
\]
i.e.\ $\X\to\D$ is a maximal unipotent (Type III) degeneration. The
monodromy-invariant space is $\ker N$, of rank $22-2=20$, and after
quotienting by the polarisation directions one obtains
$E_m$ of rank $18$; the full data is given in item (d) below, using the
dual complex $\Gamma(\X_0)$ of the central fibre:
\item $\X_0=R_1\cup_C R_2$ is the union of two rational surfaces glued along
a smooth anticanonical elliptic curve $C$ (the double curve), with
$C^2|_{R_1}+C^2|_{R_2}=0$. The dual complex $\Gamma(\X_0)$ is homeomorphic
to $S^2$: it has a cell structure with $2$ monogonal $2$-cells (one for
each component $R_1,R_2$), $1$ loop edge (the double curve $C$), and $1$
vertex (Euler characteristic $1-1+2=2$); its $1$-skeleton is a circle.
Subdividing this circle into $4$ edges and $4$ vertices turns each
monogonal face into a $4$-gon and yields the ``pillowcase'' cell
structure that we shall match with the dual complex on the fibration side
in Theorem~\ref{thm:loop}(b).
\item The filtrations of [GT, Definition 5.1(4)] for $\X\to\D$ are:
$E_m(\X)=W_2'(\X)$ of rank $18$, with subfiltration $W_2 = \bZ\, v$ where
$v\in H^2$ is the class restricting to $0$ on $\X_t$ and generating the
coniveau-$1$ part, i.e.\ the class of the double curve $C$; and
\[
E_m(\X)/W_2(\X)\;\simeq\;N
\quad\text{as a lattice},
\qquad
E_m'(\X)/E_m(\X)\;\simeq\;(-1)U,
\]
where the last isomorphism is computed from the mixed Hodge structure of
$\X_0$ via the Clemens--Schmid sequence.
\end{enumerate}
\end{proposition}

\begin{proof}
(a) The Noether--Lefschetz loci in the $H$-polarised moduli space are
the hypersurfaces where $\NS(\X_t)$ properly contains $H$; outside their
union, and hence for very general $t$ in the punctured disc
$\D\setminus\{0\}$, one has $\NS(\X_t)=H$. The very-generality hypothesis
of Theorem~\ref{thm:main} is precisely that the period map of $\X$,
extended across $0$ into the Alexeev--Brunyate--Engel compactification,
sends $0$ into a point of a Type III boundary stratum lying on no
Noether--Lefschetz locus; the first sentence of the statement follows,
and the boundary-stratum assertion is a restatement of the hypothesis.

(b) A Tyurin degeneration is by definition a Type III degeneration whose
central fibre is a union of two components meeting along a smooth double
curve (the comb model). The local invariant cycle theorem and the
Kulikov--Persson--Pinkham classification apply since $\X$ admits a
Kulikov model after finite base change; for such a model $\rk N \le
20 - \rk H$ with equality for a maximal cusp, and the cusp being
Type III is equivalent to $\rk N^2 = 1$, $N^3 = 0$.

(c) is the definition of a Tyurin degeneration combined with
Kulikov--Persson--Pinkham: after base change $\X_0$ is a simple normal
crossing divisor with dual complex a sphere $S^2$; for the very general
$\X$ the two components are smooth rational surfaces meeting transversely
along a smooth elliptic anticanonical curve $C$, and there are no triple
points. The homeomorphism type is therefore $S^2$, and the minimal cell
structure described in the statement (two monogonal cells glued along the
single loop edge) has Euler characteristic $1-1+2=2$. Since
Definition 5.1(3) requires only a \emph{homeomorphism} (plus a
checkerboard structure), we are free to subdivide; the four-edge
``pillowcase'' is a common subdivision of the cell structures from both
sides, which is what Theorem~\ref{thm:loop}(b) uses. The identity
$C^2_{R_1}+C^2_{R_2}=0$ is the mirror of the triple point formula and
follows from $(\X_s\cdot C)=0$ restricted to the normal crossing.

(d) The Clemens--Schmid exact sequence for a semistable model
$\widetilde\X\to\widetilde\D$ computes the monodromy weight filtration;
taking the quotient by the polarising $H$ (which is monodromy-invariant,
since it consists of algebraic classes constant in the family) gives
$E_m$ of rank $20-2=18$. The class $v$ of the double curve $C$ (which
equals $[\X_s]|_{\X_0}$ and generates $N^2 H^2(\X_t)^\vee$ under the
identification of [GT, Remark~4.4(ii)]) spans $W_2$. The precise
lattice-theoretic content of the two displayed isomorphisms is the
computation carried out in [GT, Example~1.2] for the very general
$H$-polarised Tyurin degeneration, which we quote as an external result
with these hypotheses: there it is shown, by the Clemens--Schmid
sequence
\[
0\longrightarrow H^2(\X_t)/W_2 \longrightarrow
H^2(R_1)\oplus H^2(R_2)\xrightarrow{\;\mathrm{res}\;}H^2(C)
\]
and the gluing of the two rational components along the anticanonical
double curve, that $E_m(\X)/W_2(\X)\simeq H\oplus E_8^{\oplus2}=N$ and
that $E_m'(\X)/E_m(\X)\simeq(-1)U$, with $W_2$ a primitive isotropic
line; the rank data $\rk E_m(\X)=18$, $\rk E_m'(\X)=20$ are the
consequence used below.
\end{proof}

\section{The allowable loop and the mirror pairing}\label{sec:loop}

\begin{theorem}\label{thm:loop}
Let $\gamma\subset\PP^1$ be a simple closed loop such that the bounded
disc $D$ it encloses contains exactly the four $I_1$-points
$p_1,\dots,p_4$ and neither $0$ nor $\infty$ (the $II^*$-points). Then:
\begin{enumerate}[label=(\alph*),leftmargin=2em]
\item $\gamma$ is allowable in the sense of [GT, Definition~4.11];
\item the dual complex $\Gamma(Y,\gamma)$ is homeomorphic to $S^2$;
\item the $N$-quasipolarisation is compatible with $\pi$ and $\gamma$ in the
sense of [GT, Definition~4.21];
\item the induced filtrations satisfy
\[
E_m(Y,\gamma)\simeq E_m(\X),\qquad
W_2(Y,\gamma)=\bZ\, f,\qquad
E_m(Y,\gamma)/W_2\simeq N,\qquad
E_m'(Y,\gamma)/E_m(Y,\gamma)\simeq (-1)U,
\]
matching Proposition~\ref{prop:deg-side}(d) under the isometry of
Proposition~\ref{prop:DN}.
\end{enumerate}
\end{theorem}

\begin{proof}
We address each item.

\textbf{(a) Allowability.} By [GT, Definition~4.11], applied with base
change the order-$6$ cover ramified at $\{0,\infty\}$ making the two $II^*$
fibres semistable ($I_0^*$ after order $6$, i.e.\ trivial monodromy) and
pulling each $I_1$ to an $I_6$, the conditions for allowability of a loop
for a fibration with section are: (i) $\gamma$ is a simple loop in the
smooth locus of $\pi$ on $\PP^1$; (ii) the monodromy around $\gamma$ is a
product of the local monodromies enclosed and, after the allowable base
change, is a maximally unipotent monodromy of Type III; (iii) the bounded
region is a disc containing in its interior only singular fibres of Kodaira
type $I_n$ ($n\ge1$) after semistable reduction (here: four points, of type
$I_6$ after base change); (iv) the total space over $\bar D$ carries a
relatively minimal d-semistable model. Properties (i) and (iii) are
immediate: $\gamma$ is smooth, bounds a disc containing only the four
$I_1$-points. For (ii), the monodromy around $\gamma$ equals the inverse of
the product of the monodromies at $0$ and $\infty$ (as
$\pi_1(\PP^1\setminus\{0,\infty,p_i\})$ is generated by loops around all
six points with one relation); each $II^*$-monodromy has order $6$ in
$SL_2(\bZ)$, its sixth power is the identity, and the product of the two is
a conjugate of a Type III unipotent after the order-$6$ base change; the
four $I_1$-monodromies are independent transvections, and their product is
maximally unipotent because they surround the complement of the two
order-$6$ points---equivalently, after the base change each becomes an
$I_6$-fibre whose monodromy is $T^6$-conjugate, and the four vanishing
cycles are in general position because $Y$ is very general, so their
intersection matrix has corank exactly $2$, giving $\rk N_\gamma =
2$, $\rk N_\gamma^2 = 1$: a Type III cusp. Property (iv) holds by the
semistable reduction theorem applied to the d-semistable surface
$\pi^{-1}(\bar D)$: the components are smooth over the disc interior, and
the central fibre of the degeneration obtained by shrinking $D$ to a point
(the Type III K3 associated to $\gamma$ in [GT, Construction~4.1]) is
d-semistable.

\textbf{(b) Dual complex.} By [GT, Construction~4.1 and Proposition~4.7],
the dual complex $\Gamma(Y,\gamma)$ of the degeneration associated to the
allowable loop is built from the components of the d-semistable model:
vertices are the irreducible components (which over the disc $D$ are the
nef model: one component mapping to the interior), and $2$-cells are
indexed by the $I_1$-fibres, i.e.\ the nodes. For a genus-one fibration
over a disc with section and $n$ fibres of type $I_1$, the dual complex is
the ``pillowcase'': two $2$-cells (the two components of the central fibre
of the associated Type III degeneration, dual to the two rulings of the
base locus), glued along their common boundary formed by the $n$ double
curves. Here $n=4$, so $\Gamma(Y,\gamma)$ is the ``pillowcase''
complex: a quadrangulation of $S^2$ whose $1$-skeleton is the equatorial
$4$-cycle, with $4$ vertices, $4$ edges, and $2$ quadrilateral faces
(Euler characteristic $4-4+2=2$). In particular
$\Gamma(Y,\gamma)$ is \emph{homeomorphic} to the dual complex
$\Gamma(\X_0)\simeq S^2$ of Proposition~\ref{prop:deg-side}(c); it is a
subdivision of the minimal cell structure of $\Gamma(\X_0)$ (whose single
loop edge is subdivided into the equatorial $4$-cycle), in the notation
of [GT, \S4.11--4.14].

\textbf{(c) Compatibility.} By [GT, Definition~4.21], compatibility of the
$H\oplus E_8^{\oplus2}$-quasipolarisation with $(\pi,\gamma)$ requires:
(i) the quasipolarisation class, restricted to the lattice $N$, has its
$H$-component a semiample class $h=\sigma+kf$ for some integer $k\ge2$;
(ii) the $E_8^{\oplus2}$-component is contracted by $h$ and is precisely
the root system of the fibres over $\PP^1\setminus D^\circ$, i.e.\ over the
unbounded component containing the two $II^*$-fibres (equivalently, the
$E_8$'s are the fibre roots outside $\gamma$---here $0,\infty\notin D$);
(iii) after identifying $N$ with the sublattice of $\NS(Y)$, the classes in
$E_m(Y,\gamma)$ orthogonal to the fibres inside $\gamma$ generate the
$E_8^{\oplus 2}$-part. We check these with $h=\sigma+2f$
(Lemma~\ref{lem:quasi}): (i) holds with $k=2$; (ii) the two $E_8$'s are
the non-identity components of the fibres over $0,\infty$, which lie outside
$D$ since $\gamma$ encloses only the $p_i$; (iii) the classes of
$\NS(Y)=N$ orthogonal to the four vanishing (nodal) classes are, since the
$I_1$-fibres have no reducible components, exactly $N$ itself generated by
the $U$-part and the fibre roots of the two $II^*$'s---the stated lattice.
Hence compatibility holds.

\textbf{(d) Filtrations.} The filtration $E_m(Y,\gamma)$ is, by
[GT, Definition 4.22 and Proposition 4.23], the lattice
\[
E_m(Y,\gamma) \;=\;
\bigl\{ \text{classes in } H^2(Y,\bZ) \text{ trivial on the fibre-generic
part inside } \gamma \bigr\}
\;=\; N,
\]
because $N$ is generated by the section $\sigma$, the fibre $f$, and the
components of the two $II^*$-fibres, all of which restrict trivially to the
smooth fibres over the loop (the vanishing cycles of the four $I_1$-fibres
are null-homologous in $H^2(Y)$; and all of $\NS(Y)=N$ is
monodromy-invariant under $\gamma$: the $\gamma$-monodromy fixes
$f$, acts trivially on the $E_8$-components of the fibres outside $D$, and
its invariants in the $I_1$-disc directions coincide with $\bZ f$). The
fibre line $W_2(Y,\gamma)=\bZ f$ is the monodromy-fixed part, and
$E_m'(Y,\gamma)$ is the extension of $E_m(Y,\gamma)$ by the section
direction with $E_m'/E_m\simeq U(-1)$, in the conventions of
[GT, Definition~5.1(4)].

\smallskip

We now verify Definition 5.1(4) on one explicit lattice. Put
\[
\cE \;\coloneqq\; N\;\oplus\;\bigl(\bZ f'\oplus\bZ g\bigr),
\qquad
(f')^2=0,\quad f'\cdot f=0,\quad f'\cdot g=0,\quad
f\cdot g=1,\quad g^2=-2,
\]
with both $f'$ and $g$ orthogonal to
$\bZ\sigma\oplus E_8^{\oplus2}$: an overlattice of $N$ of rank $20$,
i.e.\ containing $N$ as a corank-$2$ sublattice, with distinguished
isotropic line $W_2^{\cE}\coloneqq\bZ f$. In this model we set
$E_m\coloneqq N$, $W_2\coloneqq\bZ f$, $E_m'\coloneqq\cE$. The key
form computation is
\[
\cE/\bZ f \;=\; N\oplus U(-1),
\]
an \emph{orthogonal} direct sum: since $f'$ and $g$ are both orthogonal
to $\sigma$, $f$, and $E_8^{\oplus2}$, the two summands are orthogonal,
and on $\bZ f'\oplus\bZ g$ the induced form has gram matrix
$\begin{pmatrix} 0 & -1\\ -1 & -2\end{pmatrix}$
in the basis $f',g$ (the sign being that of [GT, Definition~5.1(4)],
which reverses the form); with the basis change $e_1=f'$, $e_2=g-f'$
this becomes
$\begin{pmatrix} 0 & -1\\ -1 & 0\end{pmatrix}=U(-1)$.
Note $W_2$ lies in $E_m$, so the quotient $E_m/W_2=N$ agrees with the
first summand. We embed the filtrations of both sides in
$(\cE,W_2^{\cE},f)$ as follows.
\emph{Fibration side:} send $E_m(Y,\gamma)=N$ identically and
$W_2(Y,\gamma)=\bZ f$ to $W_2^{\cE}$ (the class $f$ maps to $f$).
\emph{Degeneration side:} under the identification
$E_m(\X)/W_2(\X)\simeq N$ of Proposition~\ref{prop:deg-side}(d), map the
class $v$ of the double curve (generating $W_2(\X)$) to $f$. In
particular this matches the quoted degeneration-side data: $\rk E_m=18$,
$\rk E_m'=20$, $E_m'/E_m\simeq U(-1)$, exactly as in
Proposition~\ref{prop:deg-side}(d).

Take $\Phi\colon E_m(\X)\otimes\bQ \to E_m(Y,\gamma)$ to be the isometry
so induced, i.e.\ the identity on $N$ together with $v\mapsto f$. We
check the requirements of [GT, Definition~5.1(4)], item by item.
(i) $\Phi$ matches the distinguished isotropic classes: $v\mapsto f$ by
definition, so $W_2(\X)\xrightarrow{\sim}W_2(Y,\gamma)$. (ii) On the
quotients, $\Phi$ induces the identity of $E_m/W_2$ on both sides,
since $\cE/W_2^{\cE}=N\oplus U(-1)$ and the first factor $N=E_m/W_2$ is
mapped identically by both embeddings. (iii) The induced pairing
$E_m/W_2 \times E_m'/E_m \to \bZ$ is computed in the model on
generators: the only nonzero pairings of a class of $E_m/W_2=N$ with
a class of $E_m'/E_m=U(-1)=\bZ e_1\oplus\bZ e_2$ are
\[
f\cdot e_2 \;=\; f\cdot(g-f') \;=\; 1-0 \;=\; 1,
\qquad
f\cdot e_1 \;=\; f\cdot f' \;=\; 0,
\]
and $0$ on $\sigma$, $f$-multiples and the $E_8^{\oplus2}$-roots (since
$\sigma,f$ and all roots are orthogonal to $f'$ and $g$). Thus the
pairing has gram matrix $\begin{pmatrix}\sigma&f\end{pmatrix}
\mapsto$ rows against columns $e_1,e_2$ given by
$\bigl(\begin{smallmatrix}0&0\\0&1\end{smallmatrix}\bigr)$ on the
$\{\sigma,f\}$-part and zero on $E_8^{\oplus2}$, on \emph{both} sides:
on the fibration side by the intersection form of $Y$, and on the
degeneration side because the quoted [GT, Example~1.2] data carry the
identical intersection form (there $v$ pairs as $v\cdot v'=1$ with the
double-curve partner, matching $f\cdot e_2=1$). Hence $\Phi$ preserves
the pairing between the two graded pieces. (iv) The quasi-polarisation
classes correspond: on the fibration side the semiample class is
$h=\sigma+2f\in H\subset N$ (Lemma~\ref{lem:quasi}); on the degeneration
side the $H$-polarisation restricted to $E_m(\X)$ is the pullback of
$h=\sigma+2f$ along the isomorphism
$E_m(\X)/W_2(\X)\xrightarrow{\sim}N$ (it is the generator of the ample
ray fixed by the $H$-polarised family hypothesis); the indeterminacy of
this pullback by multiples of $v$ is exactly $W_2(\X)$, which is killed
in the quotient, and under $v\mapsto f$ both map to $0$ in
$E_m/W_2=N$. Hence $\Phi$ sends the degeneration-side polarisation to
the fibration-side semiample class. Thus $\Phi$ respects both
filtrations, both $W_2$-lines, the pairings of the quotients, and the
polarisation classes: this is exactly [GT, Definition~5.1(4)].
\end{proof}

\begin{theorem}[{Giovenzana--Thompson, [GT, Theorem C]}]\label{thm:GT}
Let $L,N$ be Dolgachev--Nikulin dual even lattices with
$N=L\oplus{}$(unimodular factors) and
$\rk N\le 18$, so that very general $L$-polarised K3 surfaces admit maximal
Type III degenerations and very general $N$-quasipolarised K3 surfaces have
$\NS=N$. If
\begin{enumerate}[label=(\roman*),leftmargin=2.2em]
\item $\X\to\D$ is a very general $L$-polarised Tyurin degeneration, with
dual complex $\Gamma(\X_0)$ and filtrations $E_m(\X)\subset E_m'(\X)$ from
the monodromy weight filtration modulo the polarisation;
\item $\pi\colon Y\to\PP^1$ is the elliptic fibration (unique up to
isomorphism, as $W(Y)=0$) on a very general $N$-quasipolarised K3 surface
in which the quasipolarising roots are contracted, with singular fibres of
types read off from the root system of $N$;
\item $\gamma\subset\PP^1$ is an allowable loop with dual complex
$\Gamma(Y,\gamma)\simeq\Gamma(\X_0)$ and with isometric filtrations
$E_m(Y,\gamma)\simeq E_m(\X)$, $W_2\leftrightarrow W_2$;
\end{enumerate}
then $\X\to\D$ and $(Y,\pi,\gamma)$ are a mirror pair in the sense of
[GT, Definition 5.1]: there is a checkerboard structure on $\Gamma$, the
quasipolarisation is compatible with $(\pi,\gamma)$, and the limiting mixed
Hodge structures correspond under the mirror diffeomorphism of the
SYZ-fibration, i.e.\ conditions (1)--(4) of Definition 5.1 hold.
\end{theorem}

\begin{remark}
This is [GT, Theorem C(iii)$\Rightarrow$(i)] together with [GT,
Definition 5.1], specialised to the case $(L,N)=(H,H\oplus E_8^{\oplus2})$.
Hypothesis (i) is Proposition~\ref{prop:deg-side}; hypothesis (ii) is
Theorem~\ref{thm:fibration}; hypothesis (iii) is
Theorem~\ref{thm:loop}(a),(b),(d); the Dolgachev--Nikulin duality
requirement of their Definition 5.1(1), namely that $N$ be
DN-dual to $L=H$ with mirror lattice $\bar N=H$, is
Proposition~\ref{prop:DN}. The checkerboard condition of their
Definition 5.1(3) is a bipartition of the vertices of the relevant
$1$-skeleton together with the light-leaf condition; here the
$1$-skeleton is the equatorial $4$-cycle, which is bipartite because $4$
is even (alternate the two colours around the four double-curve
vertices, in the notation of [GT, \S4.13--4.14]), and the light-leaf
condition is satisfied because each of the two leaf branches of
$\PP^1\setminus\gamma$ carries fibres of type $II^*$ only, whose
associated leaves are of the same type, so the condition holds by
its formulation in [GT, Definition 4.13]. Finally the compatibility
condition of their Definition 5.1(2) is Theorem~\ref{thm:loop}(c).
\end{remark}

\begin{corollary}\label{cor:final}
The very general $H$-polarised Tyurin degeneration $\X\to\D$ and the
fibration $(Y,\pi,\gamma)$ of Theorems \ref{thm:fibration} and
\ref{thm:loop} are a mirror pair.
\end{corollary}

\begin{proof}
All hypotheses of Theorem~\ref{thm:GT} are verified:
Dolgachev--Nikulin duality by Proposition~\ref{prop:DN}; the structure of
the degeneration by Proposition~\ref{prop:deg-side}; existence, uniqueness
and fibre types of the fibration by Theorem~\ref{thm:fibration} (with the
explicit Weierstrass model of Lemma~\ref{lem:weierstrass}); allowability,
dual complex, compatibility and filtration matching by
Theorem~\ref{thm:loop}. Hence the conclusion of Theorem~\ref{thm:GT}
applies and Theorem~\ref{thm:main} follows.
\end{proof}

\begin{remark}[Independence of the ``checkerboard'' and Hodge-theoretic
clauses]
For completeness we record the remaining clauses of [GT, Definition 5.1]
as they specialise here. The monodromy action on $E_m(Y,\gamma)$ computed
in Theorem~\ref{thm:loop}(d) is generated by the transvection in the fibre
class $f$ (the $N_\gamma$ of a Type III cusp with vanishing $I_6$
cycles), matching the degeneration side $N$ with $\rk N = 2$, $\rk N^2=1$;
the induced action on $E_m/W_2\simeq N$ is trivial on the $E_8^{\oplus2}$
factor and the unipotent transvection $T_f$ on the $U$-factor, and dually
on $E_m'/E_m\simeq U(-1)$. On the degeneration side the same holds with $v$
in place of $f$, by the local invariant cycle theorem. Thus the full data
of [GT, Definition 5.1(4)] coincide via our $\Phi$. The Hodge filtration
matching (their clause concerning the period point of $Y$ corresponding to
the limiting point of $\X\to\D$ under the cusp identification) is
precisely the content of the mirror map for the self-dual pair
$(H,N)$ of Proposition~\ref{prop:DN} and holds by [GT, Theorem B], whose
hypotheses are again those verified above. This completes the verification
of every clause of Definition 5.1 in this case.
\end{remark}

\section*{References}

\begin{itemize}[leftmargin=2em]
\item[{[ABE]}] V.~Alexeev, A.~Brunyate, P.~Engel,
\emph{Compactifications of moduli of elliptic K3 surfaces: Stable pair and
toroidal}, Geom. Topol. \textbf{26} (2022), 3525--3588.
\item[{[GT]}] L.~Giovenzana, A.~Thompson,
\emph{Degenerations and fibrations of K3 surfaces: Lattice polarisations and
mirror symmetry}, arXiv:2405.12009. [Theorems B and C; Definitions 4.11,
4.21, 5.1]
\item[{[Kod]}] K.~Kodaira, \emph{On compact analytic surfaces II, III},
Ann.\ of Math.\ \textbf{77} (1963), 563--626; \textbf{78} (1963), 1--40.
\item[{[KPP]}] V.~Kulikov, \emph{Degenerations of K3 surfaces and Enriques
surfaces}, Math.\ USSR-Izv.\ \textbf{11} (1977), 957--989;
U.~Persson, H.~Pinkham, \emph{Degeneration of surfaces with trivial
canonical bundle}, Ann.\ of Math.\ \textbf{113} (1981), 45--66.
\item[{[Ni]}] V.~V.~Nikulin, \emph{Integer symmetric bilinear forms and
some of their geometric applications}, Math.\ USSR-Izv.\ \textbf{14}
(1980), 103--167.
\item[{[PS--Sh]}] I.~Piatetski-Shapiro, I.~R.~Shafarevich, \emph{A Torelli
theorem for algebraic surfaces of type K3}, Math.\ USSR-Izv.\ \textbf{5}
(1971), 547--588.
\item[{[SD]}] B.~Saint-Donat, \emph{Projective models of K3 surfaces},
Amer.\ J.~Math.\ \textbf{96} (1974), 602--639.
\item[{[Sh]}] T.~Shioda, \emph{On elliptic modular surfaces},
J.~Math.\ Soc.\ Japan \textbf{24} (1972), 20--59; and the Shioda--Tate
formula.
\item[{[Ty]}] A.~N.~Tyurin, \emph{Fano versus Calabi--Yau}, The Fano
Conference, 701--734, Univ.\ Torino, Turin, 2004.
\end{itemize}

\end{document}
